% IRQROUTER chapter. Configuration-driven: vector numbers, word counts and
% hart counts come from include/defines.tex. Section order is the
% per-peripheral template enforced by python/check_intro_names.py.
\section{Overview}

The interrupt router (\peripheral{IRQROUTER}) is the chip's peripheral interrupt controller. Every deglitched peripheral interrupt level terminates here, and delivery follows the claim/complete model of a RISC-V platform-level interrupt controller: the router raises one external-interrupt wire (\texttt{meip}, vector \MeipVector{}) to each hart, and the servicing hart reads \register{CLAIM} to discover and atomically claim the pending source. Routing is per hart and per vector, so a peripheral's interrupt can be steered to whichever hart currently owns that peripheral.

\begin{itemize}
	\item One row of \IrqEnableWordsWord{} 32-bit enable words per hart (\register{HxENL}, \register{HxENM}, \register{HxENU}, \register{HxENX}), one bit per vector, covering vectors 0 to \TopVector{}
	\item One \texttt{meip} wire per hart (vector \MeipVector{}); \register{CLAIM} read claims the lowest pending enabled vector, \register{CLAIM} write completes it
	\item Exactly-once delivery: a claimed vector is hidden from every hart until completed
	\item Fixed priority: the lowest pending vector number wins
	\item Read-only raw-level words (\register{PENDL}, \register{PENDM}, \register{PENDU}, \register{PENDX}) and under-service words (\register{INSVCL}, \register{INSVCM}, \register{INSVCU}, \register{INSVCX})
	\item Resets fully masked, so the router is inert until programmed
	\item Runs on the free-running MCLK, so a routed interrupt wakes a hart whose clock is stopped by \asminline{sleep}
\end{itemize}

The \peripheral{CLINT} vectors (\ClintMsipVector{} and \ClintMtipVector{}) are never delivered through \texttt{meip}: each hart receives its own software and timer interrupts on dedicated wires, so the corresponding row bits are writable but have no effect.

\section{Block diagram}

Figure \ref{fig:irq-fabric-diagram} draws the router in its place in the chip: every interrupt source on one side, one routing row per hart and the claim/complete stage in the middle, and the harts on the other. The two \peripheral{CLINT} levels are drawn around the router rather than through it.

% Generated by LatexUserGuide.GenerateIrqFabricDiagram (make chip).
\begin{figure}[htpb]
	\centering
	\resizebox{\linewidth}{!}{\input{include/IrqFabricDiagram.tex}}
	\caption{The interrupt fabric: sources, routing rows, claim/complete stage and harts.}
	\label{fig:irq-fabric-diagram}
\end{figure}

\section{Functional description}

\subsection{Routing rows}

Hart $h$'s row is the \IrqEnableWordsWord{} words \register{HxENL} (vectors 0 to 31), \register{HxENM} (32 to 63), \register{HxENU} (64 to 95) and \register{HxENX} (96 to \TopVector{}), with $x = h$. Bit $v \bmod 32$ of the word selected by $v \div 32$ enables vector $v$ for that hart. The packing is contiguous, so bit $b$ of \register{HxENU} is vector $64 + b$. Any hart may write any row through the shared window, which is what allows a supervising hart to hand a peripheral to another hart. The rows start at the block base, 16 bytes per hart; \register{CLAIM} and the status words sit at the fixed offsets \texttt{0x800}, \texttt{0x810} and \texttt{0x820}, independent of the hart count.

\subsection{Claim and complete}

Hart $h$'s \texttt{meip} wire is asserted while some vector $v$ is pending (deglitched level high), enabled in row $h$, and not under service. A read of \register{CLAIM} by hart $h$ returns the lowest such $v$ and marks it under service, which hides it from every hart's \texttt{meip} until it is completed; the arbiter attributes the read to the requesting hart, so the router evaluates the reader's own row. A read that finds nothing pending for the reader returns \texttt{0xFFFFFFFF}, which happens when another hart claimed the only pending source first. A write of $v$ to \register{CLAIM} completes it: the vector becomes deliverable again, and if its level is still asserted it pends again immediately. Figure \ref{fig:irq-claim-complete} shows one delivery.

% Generated by LatexUserGuide.GenerateIrqClaimCompleteDiagram (make chip).
\begin{figure}[htpb]
	\centering
	\input{include/IrqClaimCompleteDiagram.tex}
	\caption{Claim/complete delivery of one peripheral interrupt.}
	\label{fig:irq-claim-complete}
\end{figure}

Completion is not qualified by owner, so that a supervising hart can complete a vector on behalf of a hart that died while servicing it; the under-service words show which vectors are stuck. Every hart, hart 0 included, uses this path; row 0 is hart 0's live routing row.

\subsection{Status words}

The four \register{PENDx} words carry the raw deglitched level of every vector, whether enabled or not, and the four \register{INSVCx} words carry the under-service state. Both sets are read-only and exist for diagnosis and recovery.

\section{Interrupts and events}

The router does not raise a vector of its own; it drives vector \MeipVector{} (\texttt{meip}) into each hart. Table \ref{t:irqrouter-irq} states the condition.

\begin{table}[htb]
	\centering
	\caption{\peripheral{IRQROUTER} delivery into the harts.}
	\label{t:irqrouter-irq}
	\begin{tabular}{ l p{4.4cm} p{3.8cm} p{4.6cm} }
	\textbf{Vector} & \textbf{Condition} & \textbf{Set by} & \textbf{Cleared by} \\ \hline
	\MeipVector{} & A vector pending, enabled in this hart's row and not under service & The peripheral asserting its interrupt level & Reading \register{CLAIM} (hides the vector until completed) \\
	\hline
	\end{tabular}
\end{table}

\section{Programming}

\subsection{Routing a peripheral to hart $h$}

\begin{enumerate}
	\item Clear the peripheral's vector bits in the previous owner's row (\register{HxENL} to \register{HxENX} with $x$ the previous owner).
	\item Set the same bits in hart $h$'s row.
	\item Enable vector \MeipVector{} in hart $h$'s interrupt controller.
\end{enumerate}

\subsection{Servicing a routed interrupt}

\begin{enumerate}
	\item Read \register{CLAIM}; if the value is \texttt{0xFFFFFFFF}, return, because the interrupt was spurious.
	\item Service the source and clear the interrupt flag at the peripheral.
	\item Write the claimed vector number to \register{CLAIM}.
\end{enumerate}

\section{Usage notes}

\paragraph{Clear the level at the peripheral before completing.} A completed vector whose level is still asserted pends again at once, so a handler that completes first re-enters immediately.

\paragraph{Treat \texttt{0xFFFFFFFF} from \register{CLAIM} as spurious.} It is the normal outcome when several harts enable the same vector and another hart claimed it first.

\paragraph{Route a vector to one hart at a time unless the sources are idempotent.} Exactly-once delivery guarantees that one hart services each event, not which hart.

\paragraph{Recover a stuck vector from a supervisor.} Read the \register{INSVCx} words to find a vector left under service by a dead hart and write its number to \register{CLAIM} from any hart.
